% Part: computability
% Chapter: recursive-functions
% Section: primitive-recursion

\documentclass[../../../include/open-logic-section]{subfiles}

\begin{document}

\olfileid{cmp}{rec}{pre}
\olsection{আদিম পুনরাবৃত্তি}

স্বাভাবিক সংখ্যার একটি বৈশিষ্ট্য হলো, উত্তরসূরি ক্রিয়া~$+1$ সসীমসংখ্যক
বার প্রয়োগ করে প্রত্যেক স্বাভাবিক সংখ্যায়~$0$ থেকে পৌঁছানো যায়—যেকোনো
স্বাভাবিক সংখ্যা হয়~$0$, নয়তো~$0$-এর কোনো উত্তরসূরির উত্তরসূরি
\dots{}। এই তথ্য ব্যবহার করে $h\colon\Nat \to \Nat$ অপেক্ষক নির্দিষ্ট
করার একটি উপায় হলো: (ক) আর্গুমেন্ট~$0$-এর জন্য $h$-এর মান নির্দিষ্ট করা,
এবং (খ) $h(x)$-এর মান দেওয়া থাকলে $h(x+1)$-এর মান কীভাবে গণনা করতে হবে,
তা নির্দিষ্ট করা। (ক) সরাসরি বলে দেয় $h(0)$ কী, তাই~$0$-এর জন্য
$h$ সংজ্ঞায়িত হয়। এখন (খ)-তে দেওয়া নির্দেশ $x=0$-এর জন্য ব্যবহার করে
$h(0)$ থেকে $h(1)=h(0+1)$ গণনা করতে পারি। একই নির্দেশ $x=1$-এর জন্য
ব্যবহার করে $h(1)$ থেকে $h(2)=h(1+1)$ গণনা করি, এবং এভাবে চলতে থাকে।
প্রত্যেক স্বাভাবিক সংখ্যা~$x$-এর জন্য শেষ পর্যন্ত এমন ধাপে পৌঁছাব, যেখানে
$h(x)$-কে $h(x+1)$ থেকে সংজ্ঞায়িত করা হয়; অতএব সব $x\in\Nat$-এর জন্য
$h(x)$ সংজ্ঞায়িত।

উদাহরণস্বরূপ, নিচের দুটি সমীকরণ দিয়ে $h\colon\Nat\to\Nat$ নির্দিষ্ট
করা হলো:
\begin{align*}
h(0) & =  1\\
h(x+1) & =  2 \cdot h(x)
\end{align*}
গুণ কীভাবে করতে হয় তা জানা থাকলে, এই সমীকরণ দুটি উপরের (ক) ও (খ)-র
জন্য প্রয়োজনীয় তথ্য দেয়। দ্বিতীয় সমীকরণটি পরপর প্রয়োগ করে পাই
\begin{align*}
  h(1) & = 2\cdot h(0) = 2,\\
  h(2) & = 2\cdot h(1) = 2\cdot 2,\\
  h(3) & = 2 \cdot h(2) = 2\cdot 2 \cdot 2,\\
  & \vdots
\end{align*}
অতএব নির্দিষ্ট করা $h$ অপেক্ষকটি হলো $h(x)=2^x$।

স্বাভাবিক সংখ্যার এই বৈশিষ্ট্য নিশ্চিত করে যে এই দুটি শর্ত পূরণকারী
$h$ একটিই আছে। এই ধরনের একজোড়া সমীকরণকে $h$ অপেক্ষকের
\emph{আদিম পুনরাবৃত্তির মাধ্যমে সংজ্ঞা} বলা হয়। নামটির কারণ হলো,
$h$-কে “পুনরাবৃত্তভাবে” সংজ্ঞায়িত করা হয়েছে—অর্থাৎ সংজ্ঞায়, বিশেষত
দ্বিতীয় সমীকরণে, ডান দিকে $h$ নিজেই আছে। আর “আদিম” বলা হয়, কারণ
$h(x+1)$ সংজ্ঞায়িত করার সময় কেবল $h(x)$-এর মান, অর্থাৎ অব্যবহিত
পূর্ববর্তী মান, ব্যবহার করা হয়। এটি~$\Nat$-এ অপেক্ষক পুনরাবৃত্তভাবে
সংজ্ঞায়িত করার সবচেয়ে সরল পদ্ধতি।

আদিম পুনরাবৃত্তির মাধ্যমে যোগ ও গুণের মতো আরও মৌলিক অপেক্ষকও
সংজ্ঞায়িত করা যায়। তবে এই ক্ষেত্রে অপেক্ষকগুলি $2$-স্থানীয়। আমরা
একটি আর্গুমেন্টের স্থান স্থির করে অন্য স্থানটিতে পুনরাবৃত্তি ব্যবহার করি।
যেমন, $\Add(x,y)$ সংজ্ঞায়িত করতে~$x$ স্থির রেখে প্রথমে $y=0$-এর জন্য
মান এবং পরে~$y+1$-এর জন্য $y$-এর সাহায্যে মান নির্ধারণ করতে পারি।
$x$ স্থির, তাই সংজ্ঞায়িত সমীকরণের বাঁ ও ডান উভয় দিকেই তা থাকবে।
\begin{align*}
\Add(x,0) & =  x\\
\Add(x,y+1) & =  \Add(x,y)+1
\end{align*}
এই সমীকরণগুলি সব $x$ এবং~$y$-এর জন্য $\Add$-এর মান নির্দিষ্ট করে।
যেমন, $\Add(2,3)$ পেতে $x=2$-এর জন্য সংজ্ঞায়িত সমীকরণগুলি প্রয়োগ
করি: প্রথমটি থেকে $\Add(2,0)=2$, তারপর দ্বিতীয়টি থেকে পরপর
$\Add(2,1)=2+1=3$, $\Add(2,2)=3+1=4$, এবং
$\Add(2,3)=4+1=5$ পাই।

$\Add$-এর সংজ্ঞায় দ্বিতীয় সমীকরণের ডান দিকে $+$ ব্যবহার করেছি, কিন্তু
শুধু~$1$ যোগ করার জন্য। অন্যভাবে বললে, আমরা উত্তরসূরি অপেক্ষক
$\Succ(z)=z+1$ ব্যবহার করে আগের মান $\Add(x,y)$-এ তা প্রয়োগ করেছি,
যাতে $\Add(x,y+1)$ সংজ্ঞায়িত হয়। তাই পুনরাবৃত্ত সংজ্ঞাটিকে এমনও
ভাবা যায় যে একটি মাত্র অপেক্ষককে আগের মানে প্রয়োগ করা হচ্ছে। তবে
অপেক্ষকটি শুধু আগের মানের উপর নয়, $x$ ও~$y$-এর উপরও নির্ভর করতে
দেওয়াতে কোনো অসুবিধা নেই—কখনও কখনও তা প্রয়োজনীয়ও। বিবেচনা করো:
\begin{align*}
  \Mult(x,0) & =  0 \\
  \Mult(x,y+1) & =  \Add(\Mult(x,y),x)
\end{align*}
$\Mult$ অপেক্ষকটির এটি একটি আদিম পুনরাবৃত্ত সংজ্ঞা: আগের মান
$\Mult(x,y)$ এবং প্রথম আর্গুমেন্ট~$x$ উভয়ের উপর $\Add$ অপেক্ষক
প্রয়োগ করা হয়েছে। এটিও সব $x$ ও~$y$ আর্গুমেন্টের জন্য $\Mult(x,y)$
সংজ্ঞায়িত করে। যেমন, $\Mult(2,3)$ নির্ধারিত হয়
$\Mult(2,0)$, $\Mult(2,1)$, $\Mult(2,2)$ ও~$\Mult(2,3)$ পরপর গণনা করে:
\begin{align*}
  \Mult(2,0) & = 0\\
  \Mult(2,1) & = \Mult(2,0+1) =
  \Add(\Mult(2,0), 2) = \Add(0, 2) = 2\\
  \Mult(2,2) & = \Mult(2,1+1) =
  \Add(\Mult(2,1), 2) = \Add(2, 2) = 4\\
  \Mult(2,3) & = \Mult(2,2+1) =
  \Add(\Mult(2,2), 2) = \Add(4, 2) = 6
\end{align*}

সাধারণ বিন্যাসটি হলো এই: $h(x_0,\dots,x_{k-1},y)$ অপেক্ষকের আদিম
পুনরাবৃত্ত সংজ্ঞা দিতে দুটি সমীকরণ দিই। প্রথমটি $h(x_0,\dots,x_{k-1},0)$-এর
মানকে $h$-এর উল্লেখ ছাড়া সংজ্ঞায়িত করে। দ্বিতীয়টি
$h(x_0,\dots,x_{k-1},y+1)$-এর মানকে $h(x_0,\dots,x_{k-1},y)$,
অন্য আর্গুমেন্টগুলি $x_0$, \dots,~$x_{k-1}$ এবং~$y$-এর সাহায্যে সংজ্ঞায়িত
করে। দ্বিতীয় সমীকরণে কেবল $h$-এর অব্যবহিত পূর্ববর্তী মান ব্যবহার করা
যায়। এই দুটি সমীকরণের ডান দিকের ক্রিয়াগুলিকে যদি নিজেরাই $f$ ও~$g$
অপেক্ষক হিসেবে ভাবি, তবে আদিম পুনরাবৃত্তির মাধ্যমে নতুন $h$ অপেক্ষক
সংজ্ঞায়নের সাধারণ বিন্যাসটি হলো:
\begin{align*}
  h(x_0, \dots, x_{k-1}, 0) & = f(x_0, \dots, x_{k-1})\\
  h(x_0, \dots, x_{k-1}, y+1) & =
  g(x_0, \dots, x_{k-1}, y, h(x_0, \dots, x_{k-1}, y))
\end{align*}
$\Add$-এর ক্ষেত্রে $k=1$ এবং $f(x_0)=x_0$ (অভিন্ন অপেক্ষক), আর
$g(x_0,y,z)=z+1$ (তৃতীয় আর্গুমেন্টের উত্তরসূরি ফেরত দেওয়া
$3$-স্থানীয় অপেক্ষক):
\begin{align*}
  \Add(x_0, 0) & = f(x_0) = x_0\\
  \Add(x_0, y+1) & = g(x_0, y, \Add(x_0, y)) =
  \Succ(\Add(x_0, y))
\end{align*}
$\Mult$-এর ক্ষেত্রে $f(x_0)=0$ (সবসময়~$0$ ফেরত দেওয়া ধ্রুবক
অপেক্ষক) এবং $g(x_0,y,z)=\Add(z,x_0)$ (শেষ ও প্রথম আর্গুমেন্টের যোগফল
ফেরত দেওয়া $3$-স্থানীয় অপেক্ষক):
\begin{align*}
  \Mult(x_0,0) & =  f(x_0) = 0 \\
  \Mult(x_0,y+1) & =  g(x_0, y, \Mult(x_0,y)) =
  \Add(\Mult(x_0,y), x_0)
\end{align*}

\end{document}
